%==============================================================================
\section{Bộ lọc phương sai thấp (Low Variance Filter)}
\label{sec:low-variance-filter}
%==============================================================================

\begin{dinhnghia}[title={Bộ lọc phương sai thấp}]
Kỹ thuật \textbf{lựa chọn đặc trưng} loại các feature có \textbf{phương sai (variance)} nhỏ hơn ngưỡng. Feature gần như \textbf{không đổi} giữa các mẫu thường ít khả năng phân biệt được lớp target.
\end{dinhnghia}

\begin{ynghia}[title={Ý nghĩa}]
Phương sai đo mức \textbf{phân tán} của một biến. Biến hằng (hoặc gần hằng) mang ít ``năng lực phân loại''; loại bỏ chúng có thể giảm nhiễu và chi phí tính toán mà không làm mất nhiều tín hiệu dự báo.
\end{ynghia}

\subsection{Thuật toán}

\begin{phuongphap}[title={Các bước triển khai}]
\begin{enumerate}
  \item Tính \textbf{phương sai} từng feature trên toàn tập $X$.
  \item Đặt \textbf{ngưỡng} phương sai (ví dụ $0{,}1$).
  \item Xác định chỉ số các feature có $\mathrm{Var} < \text{ngưỡng}$.
  \item \textbf{Xóa} các feature đó khỏi ma trận dữ liệu.
  \item Dùng tập feature còn lại để huấn luyện mô hình.
\end{enumerate}
\end{phuongphap}

\subsection{Ví dụ Python --- dataset Diabetes}

\begin{dongco}[title={Dữ liệu và ngưỡng}]
Bộ \texttt{diabetes} qua \texttt{fetch\_openml}; ngưỡng phương sai $0{,}1$.
\end{dongco}

\begin{vidu}[title={Mã đầy đủ}]
\begin{lstlisting}
import numpy as np
import pandas as pd
from sklearn.datasets import fetch_openml

data = fetch_openml("diabetes", version=1, as_frame=True, parser="auto").frame
X = data.iloc[:, :-1].values
y = data.iloc[:, -1].values

variances = np.var(X, axis=0)
threshold = 0.1

low_var_indices = np.where(variances < threshold)
X_filtered = np.delete(X, low_var_indices, axis=1)

print('Shape of the filtered dataset:', X_filtered.shape)
\end{lstlisting}
\end{vidu}

\begin{ynghia}[title={Kết quả đầu ra}]
\textbf{Output:}
\begin{lstlisting}
Shape of the filtered dataset: (768, 8)
\end{lstlisting}
Mọi feature trong diabetes OpenML đều có phương sai $\geq 0{,}1$, nên \textbf{không feature nào bị loại}.
\end{ynghia}

\subsection{Ưu và nhược điểm}

\begin{tinhchat}[title={Ưu điểm}]
\begin{itemize}
  \item \textbf{Giảm overfitting}: bỏ biến gần như không đóng góp vào dự đoán.
  \item \textbf{Tiết kiệm tài nguyên}: ít cột $\Rightarrow$ huấn luyện và suy luận nhanh hơn.
  \item \textbf{Cải thiện hiệu năng}: giảm nhiễu từ biến ``đứng yên''.
  \item \textbf{Đơn giản hóa mô hình}: tập feature gọn, dễ đọc.
  \item \textbf{Triển khai đơn giản}: chỉ cần thống kê mô tả, không cần nhãn $y$.
\end{itemize}
\end{tinhchat}

\begin{luuy}[title={Nhược điểm}]
\begin{itemize}
  \item \textbf{Mất thông tin}: biến phương sai thấp vẫn có thể quan trọng khi kết hợp phi tuyến với biến khác.
  \item \textbf{Giả định tuyến tính / độc lập}: không xét tương tác feature; loại một biến có thể phá cấu trúc ẩn.
  \item \textbf{Ảnh hưởng biến phụ thuộc}: biến ít biến thiên nhưng có quan hệ mạnh với $y$ (ví dụ mã hiếm) có thể bị xóa nhầm.
  \item \textbf{Selection bias}: ngưỡng cố định có thể không phù hợp thang đo (cần chuẩn hóa trước khi lọc).
  \item \textbf{Nhạy thang đo}: feature chưa chuẩn hóa có thể bị đánh giá sai mức phương sai tương đối.
\end{itemize}
\end{luuy}
